\chapter{Carving domains: a channel, a sphere, and a triangle mesh}
\label{ch:A6_complex_geometry}
\noindent\texttt{tutorials/A\_foundations/A6\_complex\_geometry.py}\\[0.75em]
\begin{outcome}
You can build computational domains by COMPOSITION of
signed-distance oracles (CSG booleans), choose which SIDE of a surface is
the domain (the `domain` flag), and swap in a triangulated surface (STL)
without touching the solver. Geometry in this stack is a plug-in, not a
mesh.
\end{outcome}

Three carves, same physics (the A1 Poisson MMS, order 2 each):
 1. CHANNEL: the domain INSIDE a box oracle [0.05,0.95]x[0.3,0.7] — your
    first non-square domain, and the shape of every track-D flow problem.
 2. PLATE-MINUS-SPHERE: the domain OUTSIDE a ball — the exterior
    configuration used for obstacles (note domain="outside").
 3. TRIANGLE MESH: an icosphere (a triangulated sphere) through the
    wp.Mesh/BVH backend — the same code path an STL file from CAD takes.
    For a real file: `verts, tris = your\_stl\_loader(path)` then
    `TriMeshOracle(verts, tris)`. (The classic Stanford-bunny run needs
    only that loader — see Explore (c).)

\begin{expected}
\begin{verbatim}
channel (box):        order 0.88   <- READ THIS ONE CAREFULLY
    plate minus disk:     order 2.04
    icosphere (3-D STL):  order 2.25
The smooth geometries deliver order 2. The CHANNEL does not — and that is
the lesson: a box has CORNERS, where the SDF is not differentiable, the
closest-point map is set-valued, and the Taylor shift's smoothness
assumptions fail in the corner elements. SBM is corner-degraded out of the
box (production codes add special corner treatment). Explore (b) digs in.
\end{verbatim}
\end{expected}

\section*{The script}
\begin{lstlisting}
import numpy as np

from diffsim.octree.build import build_uniform
from diffsim.mesh.nodes import build_mesh
from diffsim.mesh.constraints import build_constraints
from diffsim.mesh.basis import basis_tables
from diffsim.mesh.faces import face_tables
from diffsim.assembly.operators import DeviceMesh
from diffsim.geometry.csg import Sphere, Box
from diffsim.geometry.trimesh import icosphere
from diffsim.sbm.surrogate import classify_lambda, extract_surrogate, GeometryData
from diffsim.sbm.poisson import SBMPoisson
from diffsim.physics.poisson import l2_error_masked

DEVICE = "cuda:0"
u_star = lambda x: np.sin(np.pi * x[:, 0]) * np.cos(np.pi * x[:, 1])
f_star = lambda x: 2 * np.pi ** 2 * u_star(x)
u3 = lambda x: (np.sin(np.pi * x[:, 0]) * np.cos(np.pi * x[:, 1])
                * np.sin(np.pi * x[:, 2]))
f3 = lambda x: 3 * np.pi ** 2 * u3(x)


def carve_and_solve(oracle, level, dim, domain, u_fn, f_fn):
    tree = build_uniform(level, dim=dim)
    ret, _ = classify_lambda(tree, oracle, 0.0, domain=domain)
    sf = extract_surrogate(ret)
    mesh = build_mesh(ret, p=1)
    cons = build_constraints(mesh)
    dm = DeviceMesh.from_mesh(mesh, cons, basis_tables(1, dim=dim), DEVICE)
    geo = GeometryData.evaluate(oracle, ret, sf, face_tables(1, dim),
                                domain=domain)
    prob = SBMPoisson(dm, geo, sf, g_fn=u_fn)
    u = prob.solve(f_fn=f_fn,
                   g_outer_fn=u_fn if domain == "outside" else None)
    sgn = -1.0 if domain == "inside" else 1.0
    return l2_error_masked(dm, u, u_fn,
                           lambda x: sgn * oracle.classify(x) > 0)


if __name__ == "__main__":
    cases = [
        ("channel (inside a box)",
         Box((0.5, 0.5), (0.45, 0.2)), "inside", 2, (5, 6),
         u_star, f_star),
        ("plate minus disk (outside)",
         Sphere((0.5, 0.5), 0.25), "outside", 2, (5, 6),
         u_star, f_star),
        ("icosphere STL-path (outside, 3-D)",
         icosphere(3, (0.5, 0.5, 0.5), 0.3, device=DEVICE), "outside",
         3, (3, 4), u3, f3),
    ]
    for name, oracle, domain, dim, levels, u_fn, f_fn in cases:
        errs = [carve_and_solve(oracle, lv, dim, domain, u_fn, f_fn)
                for lv in levels]
        print(f"{name:38s}: errors "
              + "  ".join(f"{e:.3e}" for e in errs)
              + f"   order {np.log2(errs[0] / errs[1]):.2f}")
\end{lstlisting}

\begin{explore}
\begin{verbatim}
(a) Compose: channel MINUS disk (Intersection of the box-inside and the
      disk-complement — see diffsim/geometry/csg.py for Complement /
      Intersection). That domain is exactly D3's flow-past-a-cylinder
      geometry. Verify order 2 on it.
  (b) The measured channel order (0.88) is the corner effect. Verify the
      diagnosis: compute the error restricted to elements more than 3h
      from every corner (mask in l2_error_masked) — smooth order should
      reappear. Then propose a fix (corner-rounded CSG? local h-refinement
      at corners? both are used in practice).
  (c) Stanford bunny: get bun_zipper.ply / an STL (the Stanford 3D
      Scanning Repository), load vertices+triangles with e.g.
      trimesh/meshio, wrap in TriMeshOracle, scale into [0,1]^3, and rerun
      this script's case 3 with it. Report the error floor and explain it
      with the sagitta (faceting) bound.
  (d) PERFORMANCE CORNER: for case 3, time GeometryData.evaluate (the
      closest-point projections) vs level. It queries a BVH per face
      Gauss point — what exponent do you expect, and what do you measure?
\end{verbatim}
\end{explore}
